\section{Motor component}\label{sec:motor}

\subsection{Purpose}
The robot moves on two stepper motors. A stepper does not turn on its own when it receives power: its driver advances the shaft by one small step every time it receives a pulse on its STEP input, and the DIR input chooses the direction. Making a wheel turn at a given speed therefore means sending a steady train of pulses at the right rate, with DIR set the right way.

The \code{motor} component of the ESP32 firmware is the part of the code that will do exactly this, and nothing more. It knows about pins, pulses and step counts, but it does not know the size of the wheels or how the robot should move: those questions belong to a higher layer. Keeping it this simple means the same component can drive the left wheel, the right wheel, and later any other stepper on the robot.

\subsection{What has been done}
So far the work has focused on the shape of the component rather than on moving a motor:

\begin{itemize}
  \item The \textbf{public interface} has been designed and written in \code{motor.h}. It defines how the rest of the firmware describes a motor, how it creates one, and how it will later set its speed and read its position.
  \item The \textbf{creation function}, \code{motor\_new}, is written and working. It checks the description it is given, refuses anything that does not make sense, and returns a ready-to-use motor.
  \item The firmware's entry point, \code{app\_main}, now \textbf{creates the right-hand motor} at start-up, using the pin numbers from the board's pin map.
  \item The result has been \textbf{flashed and tested on the ESP32}: the start-up log confirms that the motor is created with the expected settings.
\end{itemize}

The two functions that will actually make the motor move, setting the speed and reading the position, are declared in the interface but not written yet. They come next.

\begin{center}
\begin{tabularx}{\linewidth}{@{}l l X@{}}
\toprule
\textbf{Part} & \textbf{Status} & \textbf{Notes} \\
\midrule
Public interface & \done & description of a motor, motor handle, three functions \\
Creating a motor & \done & input checks, memory allocation, tested on the board \\
Setting the speed & \todo & declared only \\
Reading the position & \todo & declared only \\
\bottomrule
\end{tabularx}
\end{center}

\subsection{The public API}
The public API is everything another part of the firmware is allowed to see and use: the contents of \code{motor.h}. Everything else stays private inside \code{motor.c}. This separation matters because it lets the inside of the component change (for example, how pulses are generated) without breaking the code that uses it.

The API is made of three ideas: a \emph{description} of a motor, a \emph{handle} that represents one motor once it exists, and three \emph{functions}.

\subsubsection{Describing a motor}
Before a motor can be created, the caller fills in a small description of it, called the motor configuration. It answers five questions:

\begin{center}
\begin{tabularx}{\linewidth}{@{}l X@{}}
\toprule
\textbf{Setting} & \textbf{What it means and why it is needed} \\
\midrule
STEP pin & The ESP32 pin wired to the driver's STEP input. The pulses will be sent on this pin. \\
DIR pin & The ESP32 pin wired to the driver's DIR input. Its level chooses the direction of rotation. \\
Steps per revolution & How many full steps the motor needs for one turn. The robot's drive motors are 1.8$^\circ$ motors, which need 200 steps; the header provides this value as a default so nobody has to remember it. \\
Microstepping & Drivers can split each full step into smaller ones for smoother motion. The value must match the setting on the driver, otherwise every distance computed from step counts will be wrong. The allowed values are 1, 2, 4, 8 and 16. \\
Invert direction & On a two-wheeled robot the motors face opposite ways, so for the robot to go forward one wheel turns clockwise and the other anticlockwise. Setting this flag on one motor hides that difference: for both motors, a positive speed then means ``forward''. \\
\bottomrule
\end{tabularx}
\end{center}

Two choices are worth noting. First, the motor is told its pins rather than looking them up itself: the pin numbers live in the board's pin map, and the code that creates the motor passes them in. The component therefore does not depend on a particular board, and can be reused elsewhere. Second, microstepping is a fixed list of named values rather than any number, so a typing mistake is caught when the code compiles or when the motor is created.

\subsubsection{The motor handle}
Once a motor has been created, the rest of the firmware refers to it through a \emph{handle}. A handle can be thought of as a ticket: the component hands it out when a motor is created, and the caller gives it back every time it wants something done to that particular motor. With two motors there are simply two handles.

The important property of the handle is that it is \emph{opaque}: the caller can keep it and pass it around, but cannot look inside it. What the motor stores internally (a copy of its configuration and the hardware resources used to generate pulses) is visible only to the motor component itself. This prevents other parts of the firmware from changing the motor's state behind its back, and it is the same pattern Espressif uses for its own drivers, so it will feel familiar to anyone who has used ESP-IDF.

\subsubsection{The three functions}
All three functions follow the same convention as the rest of ESP-IDF: they return an error code that says whether the call succeeded, and any result is written into a variable that the caller provides. The caller can therefore always check for errors, and the firmware stops with a clear message at start-up if something is misconfigured.

\begin{center}
\begin{tabularx}{\linewidth}{@{}l X@{}}
\toprule
\textbf{Function} & \textbf{Role} \\
\midrule
\code{motor\_new} & Creates a motor from a configuration and gives back its handle. It only reads the configuration and keeps its own copy, so the caller's description can be discarded afterwards. \\
\code{motor\_set\_speed} & Will set how fast the motor turns, in microsteps per second. The sign carries the direction: positive is forward, negative is backward, and zero stops the motor. \\
\code{motor\_get\_position} & Will report how far the motor has turned since it was created, counted in microsteps, going up when it moves forward and down when it moves backward. \\
\bottomrule
\end{tabularx}
\end{center}

\begin{lstlisting}[language=C]
esp_err_t motor_new(const motor_config_t *config, motor_handle_t *ret_motor);
esp_err_t motor_set_speed(motor_handle_t motor, int32_t steps_per_s);
esp_err_t motor_get_position(motor_handle_t motor, int32_t *ret_position);
\end{lstlisting}

Speeds and positions are expressed in microsteps on purpose. The motor component works at the level of the hardware; converting to metres per second, or to the robot's forward speed and turning rate, will be the job of the drive control layer above it. The position is stored as a 32-bit signed number, which at 3200 microsteps per wheel turn covers several hundred thousand turns, far more than a match will ever need.

\subsection{How a motor is created}
Creating a motor is the only part that is fully implemented, and it was written to be strict. When \code{motor\_new} is called, it works through the following steps, and stops with an error at the first problem:

\begin{enumerate}
  \item It checks that it has actually been given a configuration and a place to put the handle.
  \item It checks that the configuration makes sense: the number of steps per revolution cannot be zero, the microstepping must be one of the allowed values, and both pins must be able to act as outputs. On the ESP32, pins 34 to 39 can only be used as inputs, so a motor wired to them would never move; this check catches the mistake immediately.
  \item It reserves the memory for the new motor.
  \item It copies the configuration into that memory.
  \item Only when all of this has succeeded does it hand the new motor back to the caller.
\end{enumerate}

Each refusal is reported in the log with the reason, for example ``invalid microstep 3'', instead of a bare error code. Handing the motor back only at the very end guarantees that the caller either receives a complete, valid motor or nothing at all.

These checks were added after a real mistake during development: a configuration was once created with its microstepping left at zero, and nothing complained. With the checks in place, the same mistake now stops the firmware at start-up with a clear explanation.

\subsection{Testing}
The firmware was flashed on the ESP32 test board. At start-up, \code{app\_main} creates the right-hand motor on pins 16 (STEP) and 17 (DIR), with 200 steps per revolution and 1/16 microstepping. The start-up log confirms each stage:

\begin{lstlisting}
I (...) main: nantrobot firmware starting (status LED on GPIO 2)
I (...) motor: motor created: STEP GPIO 16, DIR GPIO 17, 200 steps/rev, 1/16 microstep
I (...) main: right motor created, handle 0x3ffb....
\end{lstlisting}


